\section{推理时引导的直觉理解}

\subsection{轨迹弯曲的几何图景}

从几何直觉上理解，推理时引导的本质是\textbf{弯曲去噪轨迹}：

\begin{enumerate}
    \item 无引导时，从 $x_T$（纯噪声）出发，扩散模型沿着边际分数 $\nabla_{x_t} \log p(x_t)$ 指示的方向，将轨迹引向数据流形上的某个点 $x_0$。
    \item 加入引导后，似然分数 $\nabla_{x_t} \log p(y|x_t)$ 在每一步给轨迹施加一个额外的``推力''，使其偏向那些满足条件 $\mathcal{A}(x_0) \approx y$ 的数据点。
    \item 最终，轨迹到达的 $x_0$ 既是数据流形上的合理点（由边际分数保证），又满足条件约束（由似然分数引导）。
\end{enumerate}

\begin{knowledgebox}{为什么扩散模型特别适合推理时引导？}
GAN 和 VAE 是一步生成模型——从噪声到数据的映射是单次前向传播。一旦这个映射固定了，就很难在推理时施加额外控制。

而扩散/流模型的生成过程是一系列\textbf{小步迭代}，每一步都可以被微调。这就像导航：一步走到目的地的话，只能走直线；但如果分很多步走，每一步都可以根据新信息调整方向——这正是推理时引导的本质。
\end{knowledgebox}

\subsection{本章小结}

推理时引导的几何本质是在去噪轨迹的每一步施加额外的方向修正力，使最终生成结果同时满足``逼真''和``符合条件''两个目标。扩散模型的迭代特性使其天然适合这种操作。


